\textbf{(CAP-COLUNI 2021 -- Q17)} A Gincana de Confraternização do CAp-COLUNI/UFV tem como objetivo a integração da comunidade escolar. As atividades são desenvolvidas por seus 480 alunos, divididos em 4 grupos, com número igual de integrantes.

Entre as provas propostas, três delas são: cantar, dançar e pular saco.

Um grupo decidiu fazer um planejamento, de acordo com a habilidade dos seus integrantes. Dessa forma, concluíram que:
\begin{itemize}
\item 84 cantam;
\item 111 pulam saco;
\item 58 dançam;
\item 53 cantam e pulam saco;
\item 36 pulam saco e dançam;
\item 36 cantam e dançam;
\item 18 cantam, dançam e pulam saco.
\end{itemize}

É CORRETO afirmar que o número de integrantes que estão aptos a praticar, exatamente, duas dessas provas é:

\begin{enumerate}
	\item 71
	\item 89
	\item 53
	\item 66
\end{enumerate}
